\documentclass[10pt,a4paper]{amsart}
\usepackage[margin=19mm]{geometry}
\usepackage{amsmath,amssymb,mathtools}
\usepackage[scr=rsfs,cal=euler]{mathalfa}
\usepackage{xfrac}
\usepackage[unicode,hidelinks,bookmarks=false]{hyperref}
\newcommand{\ie}{i.e.\ }
\newcommand{\cf}{cf.\ }
\newcommand{\resp}{respectively\ }
\newcommand{\Subsection}{\subsection}
\providecommand{\nobreakdash}{\nobreak\hbox{-}}
\setlength{\emergencystretch}{2em}
\multlinegap0pt
\usepackage{amssymb}
\usepackage[shortlabels]{enumitem}
\usepackage{mathtools}
\usepackage{dsfont}
\usepackage{xcolor}
\newtheorem{lemma}{Lemma}
\newcommand{\capacity}{\mathop{\mathrm{Cap}}\nolimits}
\title{Modulus of continuity of Monge--Amp\`ere potentials in big cohomology classes}
\author{Quang-Tuan Dang, Hoang-Son Do, Hoang Hiep Pham}
\date{Source: arXiv:2504.15763v2 (CC BY 4.0)}
\begin{document}
\maketitle

\noindent\emph{Source and licence.} Modulus of continuity of Monge--Amp\`ere potentials in big cohomology classes. Version arXiv:2504.15763v2 (CC BY 4.0), distributed by arXiv under the Creative Commons Attribution 4.0 International licence, http://creativecommons.org/licenses/by/4.0/, as recorded in the arXiv licence metadata for this exact version and shown on the abstract page https://arxiv.org/abs/2504.15763v2. Full licence notice: \texttt{licenses/arxiv-2504.15763-CC-BY-4.0.txt}.\par\smallskip

\noindent\textit{Editorial scope.} Counted unit: the article's lemma 5.3 (source0.tex lines 1058--1063). The article's complete original proof of it is delivered with it (source0.tex lines 1064--1085, one \texttt{proof} environment of 22 lines). No mathematics is written, repaired or re-derived: the statement and the proof are the article's own lines, in the article's own notation, and no retained line is substituted or re-flowed.  Retained uncounted as the source-local context the proof needs: lines 913--922; lines 1024--1038.  Nothing was omitted from any retained span, so the package carries no omission record.  No retained line contains a citation, so the wrapper carries no bibliography and no bibliography entry.  No retained line contains a figure environment, an image-inclusion command or drawing code, so no image is delivered and the package declares no asset dependency.  Editorial formatting only, in addition: an AMS \texttt{amsart} preamble that restates the article's own declarations and macros used by the retained lines, namely the environments \texttt{lemma} on separate counters, together with the standard packages those lines need.  The retained environments print in this document as Lemma 1 in order of appearance, whereas the article prints the same items with the numbering of its own full document.  The complete native source of the article as distributed in the arXiv e-print for this exact version is bundled unchanged as \texttt{sources/176897/source0.tex}.
\begin{lemma}\label{lem t^n cap}
	Assume $u, v$ are negative $\theta$-psh functions such that
	\begin{itemize}
		\item $v\leq P_{\theta}[u]$;
		\item $v=\lambda v_1+(1-\lambda) v_2$, where  $v_1, v_2$ are $\theta$-psh functions
		with $v_2$ having model singularity type and $\lambda\in (0, 1)$. 
	\end{itemize}
Then, for every $s>0$ and $0\leq t\leq 1-\lambda$, we have
	$$t^n \capacity_{v_2}\{u<v-s- t\}\leq \int_{\{u<v-s\}}\theta_u^n.$$
\end{lemma}

\begin{lemma}\label{lem sup<cap}
		Let $u, v$ be negative $\theta$-psh functions such that
	\begin{itemize}
		\item $v\leq u+M$ for some $M\geq 1$;
		\item $v=\lambda v_1+(1-\lambda) v_2$, where  $v_1, v_2$ are $\theta$-psh functions, and $v_2$ has model singularity type and $\lambda\in (0, 1/2)$;
		\item there exist $C>0$ and $\alpha>0$ such that, for every $s>0$,
		\begin{equation}\label{eq0lemsup<cap}
			\int_{\{u<v-s\}}\theta_u^n\leq C[\capacity_{v_2}(u<v-s)]^{1+\alpha}.
		\end{equation}
		\end{itemize}
	Then, for every $\varepsilon>0$,
	\begin{equation*}
		\sup_{X }(v-u)\leq \varepsilon+\dfrac{2MC^{1/n}}{1-2^{-\alpha}}[\capacity_{v_2}(u<v-\varepsilon)]^{\alpha/n}.
	\end{equation*}
\end{lemma}

\begin{lemma}\label{lem extend GZ12 prop 5.3} Under the assumption of Lemma \ref{lem sup<cap}, assume further that $v_2\geq u$. We have
the following
	inequality 
$$\sup_{X } (v_1-u)\leq \left(2+\dfrac{2 M C^{1/n}}{\lambda^{1+\alpha}(1-2^{-\alpha})}\right)\|(v_1-u)_+\|_{L^1(\nu)}^{\frac{\alpha}{(n+1)\alpha+n}},$$
where $\nu=\mathbf{1}_{\{u<v_1\}}\theta_u^n$ and $(v_1-u)_+=\max(v_1-u,0)$.
\end{lemma}

\begin{proof}
	It follows from Lemma \ref{lem sup<cap} that for every $\varepsilon>0$, we have
	\begin{equation}\label{eq: lem35.1}
	\sup_X(v-u)\leq 2\varepsilon\lambda+\dfrac{2 M C^{1/n}}{1-2^{-\alpha}}[\capacity_{v_2}(u<v-2\varepsilon\lambda)]^{\alpha/n}.
    \end{equation}
	By Lemma \ref{lem t^n cap}, we obtain
		\begin{equation}\label{eq: lem35.2}
		    \begin{split}
		        (\varepsilon\lambda)^n\capacity_{v_2}(u<v-2\varepsilon\lambda)\leq \int_{\{u<v-\varepsilon\lambda\}}\theta_u^n\leq \int_{\{u<v_1-\varepsilon\}}\theta_u^n
		    \end{split}
		\end{equation} because $u\leq v_2$ so $\{u<v-\varepsilon\lambda\}\subset\{u<v_1-\varepsilon\}$. 
   Plugging \eqref{eq: lem35.2} into \eqref{eq: lem35.1}, since $v_2\geq u$ we obtain     
        \begin{align*}
		    \sup_X(v_1-u)&\leq 
2\varepsilon+\dfrac{2 M C^{1/n}}{\varepsilon^{\alpha}\lambda^{1+\alpha} (1-2^{-\alpha})}\left(\int_{\{u<v_1-\varepsilon\}}\theta_u^n\right)^{\alpha/n}
\\
&\leq 2\varepsilon+\dfrac{2 M C^{1/n}}{\varepsilon^{\alpha}\lambda^{1+\alpha} (1-2^{-\alpha}) }\left(\frac{1}{\varepsilon}\int_{\{u<v_1-\varepsilon\}}(v_1-u)_+d\nu\right)^{\alpha/n} \\
&\leq  2\varepsilon+\dfrac{2 M C^{1/n}}{\varepsilon^{\alpha(n+1)/n}\lambda^{1+\alpha}(1-2^{-\alpha}) }\|(v_1-u)_+\|_{L^1(\nu)}^{\alpha/n},
		\end{align*} 
        where the second inequality is due to Chebyshev’s inequality.
Choosing $\varepsilon:=\|(v_1-u)_+\|_{L^1(\nu)}^{\frac{\alpha}{(n+1)\alpha+n}}$ yields the desired inequality.
\end{proof}

\end{document}
